\chapter{Introducci\'on}
\label{cap:intro}

\section{Contexto}
\label{intro:contexto}
En las últimas tres décadas, el sistema de educación superior en Chile ha experimentado un proceso de expansión y masificación, caracterizado principalmente por el crecimiento de la oferta privada. Según la \textcite{subsecretaria_de_educacion_superior_evidencias_2025},  esta expansión ha sido impulsada, en gran medida, por instituciones privadas, especialmente Centros de Formación Técnica (CFT) e Institutos Profesionales (IP). Sin embargo, tras años de crecimiento sostenido,
las instituciones enfrentan dificultades respecto a la retención y titulación oportuna: casi un tercio de los estudiantes deserta al término del primer año y la duración real de las carreras excede, en promedio, un 36\% la duración formal \parencite{subsecretaria_de_educacion_superior_evidencias_2025}.


Esta fragilidad estructural se ve hoy reforzada por una transformación del entorno económico y productivo global, impulsada por la revolución digital \parencite{araujo_chile_2023}. Dicha evolución tecnológica no solo exige una adaptabilidad
ágil de las organizaciones, sino que ha redefinido los perfiles de egreso
exigidos. En este escenario, emerge una brecha sistémica persistente: mientras la academia ha priorizado históricamente los fundamentos teóricos, la industria demanda cada vez más competencias prácticas y capacidades analíticas aplicadas \parencite{constantinov_preliminary_2015,miedema_data_2025}.

Para responder a estas nuevas exigencias, la transición hacia modelos de aprendizaje activo es respaldada tanto por lineamientos internacionales como por las demandas de los estudiantes. Por un lado, las directrices curriculares actualizadas en ciencias de la computación recomiendan adoptar enfoques de aprendizaje que enfaticen las aplicaciones prácticas y los proyectos interdisciplinarios \parencite{miedema_data_2025}.

Por otro lado, en el contexto chileno, los análisis de inserción laboral revelan que los egresados exigen una mayor \enquote{conexión teoría-práctica a través de proyectos reales} \parencite[p.~62]{latuz2024} durante
la carrera , argumentando que simular el mundo laboral y aplicar conocimientos en contextos prácticos los prepara significativamente mejor para su futuro
profesional.

\section{Motivaci\'on}
\label{intro:motivacion}
A pesar de la necesidad por alinear la educación con la industria, el proceso de mapeo curricular ha sido tradicionalmente manual, laborioso y propenso a inconsistencias \parencite{cristea_artificial_2025,russell_e_mapping_2024}.
La principal brecha actual radica en que no existen herramientas automatizadas que permitan una comparación objetiva de los programas de estudio frente a demandas laborales o estándares internacionales.
Las soluciones existentes suelen limitarse a análisis estáticos o encuestas que no logran escalar ni adaptarse a la velocidad de los cambios tecnológicos \parencite{russell_e_mapping_2024}.

Por lo tanto, resulta pertinente contar con una metodología estructurada orientada a la extracción continua de información no estructurada proveniente de datos públicos, como las mallas curriculares, lo que facilitaría su análisis y contraste de manera sistemática y periódica.
\section{Proposito de la soluci\'on}
\label{intro:proposito}
Se busca desarrollar una metodología que apoye el proceso de mapeo y análisis curricular mediante un enfoque sistemático para la identificación y comparación de competencias entre los programas académicos y las necesidades del entorno laboral.

De esta manera, se espera facilitar la labor de docentes, coordinadores y gestores académicos al proporcionar un mecanismo objetivo, consistente y dinámico para la evaluación de competencias. Asimismo, la metodología permitirá a las instituciones mantener sus programas formativos alineados con las necesidades del mercado laboral, contribuyendo a la actualización continua de sus planes de estudio y fortaleciendo la pertinencia de la formación ofrecida.